% Dieses TeX-Dokument ist utf-8 codiert. (Linux)
\documentclass[11pt,titlepage,a4,leqno] {scrartcl}
% 11pt charactergröße
% a4 mit Titelseite
% gleichungsnummern ungewöhnlicherweise links (alternativ reqno)

% in packages.tex sind die benutzten Pakete zusammengefasst
%definiere packages, die im Protokoll benutzt werden
\usepackage{ifpdf}    
\usepackage[T1]{fontenc}          % nutze LaTeX Fonts in westeuropäischer Codierung   
\usepackage[utf8]{inputenc}       % setzte Input Codierung auf utf8 - Text muss utf8 kodiert sein!                

% Mathe-Pakete
\usepackage{amsmath}
\usepackage{amssymb}

\usepackage[ngerman]{babel}       % Deutsch, neue Trennung !

\usepackage{color}                % es kann bunt werden 
\usepackage{graphicx}             % Bilder einbinden ermöglichen

\usepackage{cite}                 % Besser zitieren: -> [1-3, 6-8]

%\usepackage[textsize=small, textwidth=0.8in]{todonotes} % \todo noch zu tun

\usepackage{multirow}             % Zusammenfassen von Tabellenzellen 
                                  % \multicolumn und \multirow
\usepackage{upgreek}              % nicht-kursive Griechische Buchstaben
\usepackage{textcomp}             % \textperthousand  (promille)
\usepackage{rotating}             % rotierte Information

% das ist nur der Schönheit halber für die Bildunter-/Tabellenüberschriften !
\usepackage[format=plain,indention=1cm,scriptsize,font=sf ,labelfont=bf,
  nooneline]{caption} 
                                  % modif Bildunterschriften 
\renewcommand{\captionfont}{\small \sffamily \slshape} 
\renewcommand{\captionlabelfont}{\small \sffamily \slshape \bfseries   }

% Kontrolle der Seitenränder:
\usepackage[inner=3.5cm, outer=3cm, top=3cm, bottom=2cm, includehead, includefoot]{geometry}
\clubpenalty=10000   % letzte Zeile einer Seite nicht erste Zeile eines Absatzes
\widowpenalty=10000  % letzte Zeile eines Absatzes nicht erste Zeile einer Seite

% für Hyperlinks
\usepackage{hyperref}                            %macht anklickbare Links
\definecolor{LinkColor}{rgb}{0,0,0.75}           %dunkelblau
\hypersetup{
	pdftitle = {LaTeX-Vorlage},
	pdfsubject = {LaTeX-Vorlage},
	pdfkeywords = {LaTeX, Vorlage, Uni Wuppertal},
	pdfauthor = {Ich},
	bookmarksnumbered = true,
	bookmarksopen = false,
	colorlinks = false,
	hypertexnames = true,
	a4paper = true,
	colorlinks=true,%
	linkcolor=LinkColor,%
	citecolor=LinkColor,%
	filecolor=LinkColor,%
	menucolor=LinkColor,%
	pagecolor=LinkColor,%
	urlcolor=LinkColor,
	breaklinks=true	
}
%######################################################################
% definiere NEUE Kommandos füŸr die standard Titelseite
% Details dazu finden sich in titlepage.tex 
\newcommand{\titel} {Mein erster FP-Versuch}
\newcommand{\autoren}{Walther Stern und Otto Gerlach}
\newcommand{\betreuer}{Helmut Merkel M.Sc., Dr. Angela Kohl }
\newcommand{\datum}{\today} %datum von heute
% hier steht der Text Ihres Abstracts
\newcommand{\abstracttext}
{ 
   ~\\
   Ein Abstract ist eine gek\"urzte pr\"azise Darstellung des Inhalts einer Veröffentlichung. 
   Insbesondere werden hier die wichtigsten Resultate angegeben.\\
   Die Erdbeschleunigung wurde gemessen zu:
   % hier ist ein mathematischer Ausdruck
   \center{$g=9,81\pm 0,03_{stat.}\pm 0,02_{syst.}{\mathrm{\frac{m}{s^2}}}$}
}
\begin{document}                                       
%\pagestyle{empty} %damit kann man Seitennumerierung ausschalten
% über \input{filename} kann man einen File einbinden.
% Das wir hier nur genutzt um Das Dokument übersichtlicher zu machen.
% spezielle Titelseite füŸr das FP
% Diese Datei erzeugt die typische Uni-Wuppertal Titelseite bestehend 
% aus dem gelben Streifen mit dem Bergischen Löwen,
% dem Schriftzug "Fakultät ..." sowie Namen und Titeln



\begin {titlepage}

\unitlength1mm
\begin{picture}(150,205)(0,40)


\put(0,150){\includegraphics[height=10cm]{boxleo}}

% Trick: 5cm hohe Box mit linken umteren Punkt 5cm unterhalb 
% von Oberkante des Streifens => Oberkante text = Oberkante Balken

% Unabhängigigkeit von UTF8 oder LATIN: Universit"at
\put(25,200){\parbox[b][5cm][t]{20cm}{\textsc{\large 
Bergische Universit\"at Wuppertal \\
Fakultät für \\
{\it Mathematik und  Naturwissenschaften} \\ 
Fachgruppe Physik \\ 
~\\
Fortgeschrittenen Praktikum}}}




\put (25,185.0){\parbox[b]{20cm}{ \Large \bf \titel}}



\put (25,152) {\parbox[b]{20cm}{\large \bf \autoren}}

\put (25,125) {\parbox[t]{11cm}{\large \bf Abstract (Kurzbeschreibung) \\
                                \normalsize\rm
                                \abstracttext}}

\put (25,60) {Abgabedatum: \datum}
\put (25,65) {Betreuer: \betreuer}



\put (22,30)
{
%  \fbox{
%    \begin{turn}{-90}
    \begin{tabular}{|l|c|l|c|}
    \hline
    Bewertung Protokoll   & max. \% & +/0/- &  erreicht  \% \\  \hline \hline
    Formales              & 6      &       & \\ \hline
    Einleitung \& Theorie &  6     &       & \\  \hline        
    Durchf\"uhrung        &        &       & \\
    Auswertung		      & 33     &       & \\ 
    phys. Diskussion      &        &       & \\ 
    Zusammenfassung       &        &       & \\ \hline \hline
    Qualit"at der Messung & 15     &       & \\ \hline
    \end{tabular}
    
%    \end{turn}
%  }
}

\end {picture}
\end {titlepage}










% wenn ein Inhalts-, Bilder-, Tabellenverzeichnis gewünscht ist
% inhaltsverzeichnis
\tableofcontents
% Liste der Bilder
\listoffigures
%Liste der Tabellen
\listoftables
%
\clearpage 
%Neue Seite !
%Sorgt dafüŸr, dass alle ausstehenden Grafiken und Tabellen 
%jetzt eingefüŸgt werden, um auf einer neuen Seite ohne "Altlasten" %fortzufahren

% wie macht man Kapitel und Unterkapitel?
\section{Einleitung}\label{sec:einleitung}
Hier steht eine knappe Einleitung in die Thematik und die wesentliche
Theorie/Formeln.
Auch erste Verweise auf grundlegende Literatur~\cite{SternGerlach}.
%\autoref{SternGerlach}
Vergleichen Sie auch~\cite{wissarbeit}

\subsection{Ein UnterKapitel} \label{subsec:unterkap}
(Es ist \"ubrigens guter Ton, zwischen solche \"Uberschriften immer etwas Text
zu setzen)
\subsubsection{Ein UnterUnterKapitel}
(Alles nur zum Test)
\subsubsection{Noch ein UnterUnterKapitel}
\label{subsubsec:unterunterkap}

Was hatten wir in \autoref{subsec:unterkap} noch diskutiert?

% wie bindet man Bilder ein
\section{Experiment und Durchf\"uhrung}
Hier steht die Dokumentation des experimentellen Aubaus und evtl. wichtige 
Bemerkungen zur Durchf\"uhrung der Messung.

Statt dessen stehen hier Beispiele zu Grafiken.

Wichtig ist, das Bilder (oder auch Tabellen) in Latex nicht an einer festen Stelle stehen sollten, sondern sich bewegen \glqq floaten\grqq\ können, um
% \glqq \grqq macht "deutsche" Anführungszeichen  
ein besseres Layout zu erreichen. Man bezeichnet Figuren und Bilder deshalb auch als Floats.

% Wir definieren eine Bildumgebung
% [ tbhp ] sagt, das das Bild bevorzugt oben (t), unten (b), hier (h)
% oder auf einer neuen Seite angeordnet werden soll. Die Reihenfolge 
% ist maßgeblich
\begin{figure}[tbp]
% damit zentrieren wir das Bild
 \centering
% jetzt fügen wir ein Bild ein
 \includegraphics[angle=270, width=5cm]{sinus}
% Die Dateiendung fehlt, es wird automatisch das eps oder pdf gewählt
 \caption[Kurz: Eine Sinuskurve]
 {Lang: Dieses Bild zeigt eine Sinuskurve mit 
drei periodischen Schwingungen. Es ist 5\,cm breit. Das Bild musste um 
270\textdegree  gedreht werden. Es müsste breiter sein!}
% in eckigen Klammern steht eine Kurzbezeichnung, wie sie in einem 
% Abbildungsverzeichnis stehen köšnnte
% Wichtig ist, dass \caption eine Nummer füŸr das Bild erstellt, 
% erst dann kann man ein label darauf setzen!
\label{fig:sinus}
\end{figure}

\begin{figure}[bh]
 \centering
 \includegraphics[angle=255, width=0.25\linewidth]{sinus}
  \caption{Dieses Bild misst 25\%  der Textbreite und ist um 200\textdegree gedreht.}
 \label{fig:sinus2}
\end{figure}
Bilder haben \textbf{grundsätzlich} Unterschriften!
Schlecht ist, dass dieses Bild keine Achsenbeschriftung hat!

% wie macht man Tabellen

im Gegensatz zu Bildern haben Tabellen \textbf{grundsätzlich} Überschriften!
\begin{table}[b]
 %Auch hier unterschiedliche Texte für Überschrift und Verzeichnis:
  \caption[Die erste Tabelle hier]{Eine Tabelle mit linksb\"undiger, 
zentrierter und rechtsb\"undiger Spalte. Sie steht wegen des $[b]$ (von bottom) unten auf der Seite.
%Sie steht als erstes auf dieser 
%Seite, obwohl sie im Quelltext nach \autoref{fig:sinus} und \autoref{fig:sinus2} steht.
}
  \centering
  \begin{tabular}{|l||c|r|}
  \hline
    A & B & $\Gamma$\\
    \hline
    \hline
    \hline
    ddd & eee & $\pi$\\
    \hline
    $\Pi$ & h & iiii\\
  \hline 
  \end{tabular}
\end{table}
\begin{table}[t]
\caption[Zwei Tabellen nebeneinander]{Hier sind 2 Tabellen mit Minipages nebeneinander gesetzt. Wie die Minipages intern/extern stehen kann man 
setzen. Die Breiten kann man per Hand setzten oder automatisch.\\ In der linken Tabelle sind Zahlen sehr gut vergleichbar sind. Der Befehl \mbox{@\{\}} unterdrückt hier Zwischenräume. Diese Tabelle ist zentriert.\\
In der rechten Tabelle ist eine Zeilenhöhe künstlich, durch einen unsichtbaren Maßstab verändert. Hier sind in einer Zeile auch Spalten zusammengefasst.}
\begin{minipage}{3cm}
% die Breite der Minipage muss spezifiziert werden
\centering
\begin{tabular}{r@{}c@{}l@{}}
\hline 100 & , & 1 \\ 
\hline 10 & , & 01 \\ 
\hline 1 & , & 001 \\ 
\hline 
\end{tabular} 
\end{minipage}
% diese Befehler erlauben horizontales schieben
% \hfill
\hspace*{3cm}
\begin{minipage}[t][3cm][r]{0.5\textwidth}
% in den eckigen klammern stehen optional 
% äußere Position, Höhe, innere Position
%\flushleft/right erlaubt auch Positionsbestimmung
\begin{tabular}{|l|l|l|l|r|}
\hline
%rule{Breite}{Höhe} setzt einen nicht sichtbaren Maßstab (Breite 0) ein
\rule{0cm}{3cm}  1& 2& 3& 4& 5 \\\hline
  Spalte 1 & \multicolumn{3}{|c|}{|Spalte 2--4} & Spalte 5\\
 \hline
\end{tabular}
\end{minipage}
\end{table}

% das wichtigste zu Formeln
\section{Formeln}\label{sec:formeln}
Für Formeln sollte man unbedingt \AmS\TeX\  verwenden, das ist einfach zu bedienen und liefert bessere Ergebnisse als Standard-\LaTeX.

Dies macht der Befehl \textsf{align}.
% Achtung hier dürfen keine Lehrzeilen eingeführt werden !
\begin{align}
 \int_V\rho\,dV&=\sum_{i=1}^n m_i \quad=\quad M  \\
\intertext{Tiefer stellt man durch ein vorgestelltes underline (\_), höher durch Circumflex \^{}. Mehrere 
Buchstaben muss man in eine geschweifte Klammer einbinden.}
  S &= k\cdot\log{W} \\
\intertext{Vor Funktionen steht immer ein Backslash, z.B. \textsf{\textbackslash exp, \textbackslash sin}:}  
 c^2&=a^2+b^2 \nonumber \\
\intertext{Oben ist mit \textsf{\textbackslash nonumber} die Numerierung unterdrückt. }
 \chi&=5 \\
  x &= y &= z \label{zeile4}\\
  z &= a &= 1 \label{zeile5}
\end{align}
In Zeile \ref{zeile4} und \ref{zeile5} stehen mehrere Gleichheitszeichen übereinander.
Das geht schöner mit einem \textsf{alignat*} ohne Gleichngsnummern (macht der *).
\begin{xalignat*}{2}
 A &= B   &\qquad  B &= C \\
 C &= D   &\qquad  D &= E 
\end{xalignat*}
%\quad und \qquad machen Abstände
Die folgenden Gleichungen sind mit  \textsf{gather} zentriert untereinander geschrieben:
\begin{gather}
 A = B   \\ 
 A+B  = D
\end{gather}
Hiermit kann man mehrere Gleichungen übereinander mit nur einer Gleichungsnummer schreiben
\begin{equation}
  \begin{split}
 A &= B   \\ 
   &= T
  \end{split}\label{eq:equation}
\end{equation}
Wenn man mal eine wirklich lange Gleichung hat benutzt man \textsf{multline}.
\begin{multline}
\text{Dies ist eine}\\
\text{wirklich\qquad\qquad}\\
\text{unheimlich}\\
\text{\qquad\qquad l\"angliche}\\
\text{Gleichung}
\end{multline}



%Der Rest steht für sich selbst:

\section{Experiment und Durchf\"uhrung}

\section{Auswertung}

\section{Diskussion}

\section{Zusammenfassung}
%das drückt den Rest nach unten auf die Seite
\vfill

% Hier ist die zitierte Literatur zusammengefasst.
% Die Reihenfolge entspricht der Reihenfolge in thebibliography
\begin{thebibliography}{12345}
%-------------------------------------------------------
\bibitem{SternGerlach} 
Walther Gerlach und Otto Stern, Zeit. f\"ur Physik A, Vol9, Nr. 1, Dezember
1922.
\url{http://www.springerlink.com/content/p72218361287275g/}
\bibitem{wissarbeit}
\url{http://de.wikipedia.org/wiki/Wissenschaftliche_Publikation}
\end{thebibliography}


\end{document}


